\documentclass[11pt]{article}
\pdfinfoomitdate=1
\pdftrailerid{}
\pdfsuppressptexinfo=15
\usepackage[T1]{fontenc}
\usepackage{lmodern,fullpage,amsmath,amssymb,amsthm}
\usepackage[title]{appendix}
\usepackage[colorlinks=true,linkcolor=blue,citecolor=blue,urlcolor=blue]{hyperref}
\hypersetup{pdftitle={Arbitrarily long avoidance of 3 in Euclid--Mullin sequences},
 pdfauthor={P. Graham},pdfsubject={Prime-seeded least-prime-factor sequences},
 pdfkeywords={Euclid-Mullin sequence, Donkin question, Chinese remainder theorem, Lean}}
\newtheorem{theorem}{Theorem}
\newtheorem{lemma}[theorem]{Lemma}
\newcommand{\lpf}{\operatorname{lpf}}
\title{Arbitrarily long avoidance of 3 in Euclid--Mullin sequences}
\author{P. Graham\\ \small\url{https://orcid.org/0009-0007-8493-1505}}
\date{October 26, 2026}
\begin{document}
\maketitle
\begin{abstract}
For a prime seed $p_1$, define $p_{k+1}$ to be the least prime divisor of
$1+p_1\cdots p_k$. For every positive integer $n$, there is a prime seed
for which none of the first $n$ terms is $3$. This answers the first of
Donkin's two questions in Cameron's problem list. The proof chooses a
finite increasing sequence of auxiliary primes and forces the sequence
to begin with the seed, $2$, and those primes. Booker and Irvine used the
Chinese remainder theorem to prescribe prime-divisor paths, including
constructions from prime starting values. Here congruences also exclude every smaller
prime factor, making the prescribed divisors the least factors.
Dirichlet's theorem then supplies a prime seed. The argument, including
the congruence construction and its translation into the sequence, is
formalized in Lean. The separate question of a single prime seed whose
sequence avoids $3$ forever remains open.
\end{abstract}

\section{Introduction}

The Euclid--Mullin rule extends a finite list of primes by the least
prime divisor of one plus their product. Starting at $5$, for example,
the sequence begins $5,2,11,3$. Mullin's original least-factor sequence
starts at $2$, so its second term is $3$ \cite{mullin}. Write $\lpf(a)$ for the least prime
divisor of an integer $a>1$. Throughout this paper a prime seed $p_1$
determines the sequence
\begin{equation}\label{eq:sequence}
 p_{k+1}=\lpf(1+p_1\cdots p_k)\qquad(k\geq1).
\end{equation}

\begin{theorem}\label{thm:main}
For every positive integer $n$, there is a prime $p_1$ for which the
sequence \eqref{eq:sequence} satisfies $p_j\ne3$ for $1\leq j\leq n$.
\end{theorem}

This is the first question of Problem~19 in Cameron's archived list,
attributed there to Steve Donkin \cite[Problem~19]{cameron}.
The seed is included among the first $n$ terms and may depend on $n$.
Cameron's later Problem~109 and blog discussion report John Bray's computations of the first
position at which $3$ appears for prime seeds up to two million
\cite[Problem~109]{cameron}, \cite{cameron-blog}.
The least seeds whose first $3$ appears at each position from $1$ to $15$
were tabulated by Labos Elemer and extended by David Wasserman
\cite{oeis}. The second question in Problem~19 asks for one prime
seed avoiding $3$ forever; it remains open.

The closest method is Booker and Irvine's use of the Chinese remainder
theorem to prescribe successive prime divisors from a chosen starting
value. Their Lemma~2.2 gives an arithmetic progression, coprime to its
modulus, of starting values for which the specified primes divide the
numbers formed by adding one to the current product, allowing any prime
factor, not necessarily the least, at each step. Remark~2.16 obtains prime
starting values for loops (two distinct paths between the same two nodes) in their graph
\cite{booker-irvine}. Here we also exclude every divisor smaller than the prescribed
prime. Counting supplies a residue for each sufficiently large prime modulus;
the essential choice is to make all auxiliary primes congruent to $1$
modulo the small primes, where counting alone does not suffice.

Paris's related Lean formalization proves that a fixed sequence of prime
factors, each at most $y$, is preserved when the starting values agree
modulo an integer divisible by every prime at most $y$
\cite[\S4.3, item~(i)]{paris}.

Other Euclid--Mullin variants use different selection rules.
For the rule permitting prime factors of sums of complementary
subproducts, Booker's Theorem~1 shows that any finite set of primes can
be extended to a sequence containing every prime \cite{booker-variant}.
Booker proved that the largest-factor sequence starting at $2$ omits
infinitely many primes \cite[Theorem~1]{booker-second}; Pollack and Trevi\~no gave
an elementary proof
\cite[\S1, Theorem (Booker)]{pollack-trevino}.

The proof chooses increasing auxiliary primes with the small-prime
congruences just described. It then selects compatible seed residues at
each prime, combines them by the Chinese remainder theorem, and obtains
a prime seed by Dirichlet's theorem. Section~2 gives the construction;
Section~3 identifies the formal theorem and its accompanying proof.

This work has not yet been reviewed by experts.
The proofs, the Lean formalization and the text were produced by AI models
(from OpenAI and Anthropic), working under the author's direction.

\section{Proof of the main theorem}

We prescribe an initial segment $p,2,q_1,\ldots,q_m$ and find a prime
seed $p$ by congruences. Fix a positive integer $m$. Let
\[
 m+2<q_1<\cdots<q_m
\]
be primes satisfying $q_i\equiv1\pmod\ell$ for every prime
$\ell\leq m+2$. Such primes exist by Dirichlet's theorem in the
progression $1$ modulo the product of those small primes. Define
\[
 C_i=2\prod_{j<i}q_j\qquad(1\leq i\leq m),
\]
where an empty product is $1$. If a seed $p$ is followed by $2,q_1,\ldots$,
the integer whose least factor must be $q_i$ is $pC_i+1$.

\begin{lemma}\label{lem:local}
For each prime $\ell$, one can choose a nonzero residue $r_\ell$
modulo $\ell$ so that $r_2=1$ and, for every $1\leq i\leq m$,
\begin{align*}
 C_i r_\ell+1&\equiv0\pmod\ell &&\text{if }\ell=q_i,\\
 C_i r_\ell+1&\not\equiv0\pmod\ell &&\text{if }\ell<q_i.
\end{align*}
\end{lemma}

\begin{proof}
There are three cases. If $\ell\leq m+2$, then $C_i\equiv2\pmod\ell$
for every $i$. For $\ell=2$, take $r_2=1$; every $C_i r_2+1$ is odd.
For odd $\ell\leq m+2$, take $r_\ell=-1$; then
$C_i r_\ell+1\equiv-1\pmod\ell$. None of the $q_i$ equals such an
$\ell$, so these choices give all the required conditions.

Suppose next that $\ell=q_k$ for some $k$. The number $C_k$ is
invertible modulo $\ell$, since its prime factors are $2$ and the
earlier, smaller $q_j$. Take $r_\ell=-C_k^{-1}$.
This residue is nonzero and makes $C_k r_\ell+1$ vanish.
For $i>k$, the product $C_i$ contains $q_k$, so
$C_i r_\ell+1\equiv1\pmod\ell$. For $i<k$, the inequality
$\ell<q_i$ does not hold, and no exclusion is required.

Finally, suppose that $\ell>m+2$ and that $\ell$ is not one of the
$q_i$. For each $i$, the equation $C_i r+1=0$ in the field
$\mathbb Z/\ell\mathbb Z$ excludes at most one residue. Together with $r=0$, at most $m+1$ residues are forbidden.
Since $m+1<\ell$, a permitted residue exists. It satisfies the
stronger condition $C_i r+1\ne0$ for every $i$.
\end{proof}

\begin{lemma}\label{lem:lift}
For primes $q_1,\ldots,q_m$ with $m+2<q_1<\cdots<q_m$ and
$q_i\equiv1\pmod\ell$ for every prime $\ell\leq m+2$, there is a prime $p>3$ whose sequence
\eqref{eq:sequence} begins $p,2,q_1,\ldots,q_m$.
\end{lemma}

\begin{proof}
Set $B=q_m$ and let $M$ be the product of all primes at most
$B$. For each prime $\ell\leq B$, choose the residue supplied by
Lemma~\ref{lem:local}. The Chinese remainder theorem gives a residue
$a$ modulo $M$ satisfying $a\equiv r_\ell\pmod\ell$ for every such
$\ell$. Every $r_\ell$ is nonzero, so $\gcd(a,M)=1$.
Dirichlet's theorem therefore gives a prime $p>3$ with
$p\equiv a\pmod M$.

The first term after $p$ is $2$, because $p$ is odd. For each $i$,
the congruences give $q_i\mid pC_i+1$ and exclude every prime
$\ell<q_i$ from dividing $pC_i+1$. All these primes were included
among the moduli, since $q_i\leq B$. Thus
\[
 \lpf(pC_i+1)=q_i.
\]
Induction on $i$ now proves the claimed prefix: after $p,2,q_1,\ldots,q_{i-1}$,
the product of the terms already obtained is exactly $pC_i$.
\end{proof}

\begin{proof}[Proof of Theorem~\ref{thm:main}]
Given $n\geq1$, choose the auxiliary primes with $m=n$ and apply
Lemma~\ref{lem:lift}. The resulting prefix is $p,2,q_1,\ldots,q_n$.
Here $p>3$, and each $q_i>n+2\geq3$.
Consequently all these terms avoid $3$, in particular the first $n$.
\end{proof}

\section{Formal verification}

The argument is formalized in Lean~4 \cite{lean} using mathlib
\cite{mathlib}. The accompanying Lean project is supplied as
\href{https://github.com/6/clank-lab/tree/cameron-donkin-v1/results/cameron-donkin}{supplementary material},
with checking instructions in \texttt{lean/README.md}.\footnote{\url{https://github.com/6/clank-lab/tree/cameron-donkin-v1/results/cameron-donkin}}
The formal theorem \texttt{Donkin.arbitraryDelay} states
\[
 \begin{aligned}
 &\forall n\in\mathbb N,\quad n>0\ \Longrightarrow\\
 &\qquad \exists p\in\mathbb N,\quad p\text{ is prime}\ \land
 \bigl(\forall j\in\mathbb N,\ j<n\Longrightarrow\mathrm{term}(p,j)\ne3\bigr).
 \end{aligned}
\]
The formal least-factor function is \texttt{Nat.minFac}.
The definition has $\mathrm{term}(p,0)=p$ and
\[
 \mathrm{term}(p,k+1)=\lpf(P(p,k)+1),\qquad
 P(p,0)=p,\quad
 P(p,k+1)=P(p,k)\lpf(P(p,k)+1).
\]
Thus $P(p,k)$ is the product of terms with indices $0,\ldots,k$;
the term with index $j$ is $p_{j+1}$ in \eqref{eq:sequence}.
For a prime seed, $P(p,k)+1\geq3$, within the domain of $\lpf$.

Lemmas~\ref{lem:local} and~\ref{lem:lift}, including the induction
establishing the prefix, are checked in Lean. Dirichlet's theorem is
supplied by mathlib's formal proof. Every mathematical step is formalized;
the proofs use only Lean's standard classical axioms, specified in the
supplementary checking notes.

\begin{thebibliography}{99}
\raggedright
\bibitem{mullin}
A.~A. Mullin, \emph{Recursive function theory (a modern look at a
Euclidean idea)}, Bulletin of the American Mathematical Society
\textbf{69} (1963), 737.
\href{https://doi.org/10.1090/S0002-9904-1963-11017-4}{doi:10.1090/S0002-9904-1963-11017-4}.

\bibitem{cameron}
P.~J. Cameron, \emph{Problems}, archived problem list, page dated August~31, 2013.
\url{https://cameroncounts.github.io/web/QM/oldprob.html}.
Accessed October~26, 2026.

\bibitem{cameron-blog}
P.~J. Cameron, \emph{A prime problem}, October~5, 2012.
\url{https://cameroncounts.wordpress.com/2012/10/05/a-prime-problem/}.

\bibitem{oeis}
Labos Elemer, \emph{A093780}, The On-Line Encyclopedia of Integer
Sequences, May~3, 2004; extended by David Wasserman, April~12, 2007.
\url{https://oeis.org/A093780}.
Accessed October~26, 2026.

\bibitem{booker-irvine}
A.~R. Booker and S.~A. Irvine,
\emph{The Euclid--Mullin graph},
Journal of Number Theory \textbf{165} (2016), 30--57.
\href{https://doi.org/10.1016/j.jnt.2016.01.013}{doi:10.1016/j.jnt.2016.01.013}.
Numbering follows
\href{https://arxiv.org/abs/1508.03039v2}{arXiv:1508.03039v2}.

\bibitem{paris}
M.~Paris, \emph{The Euclid--Mullin sequence as an orbit: reformulations,
a composite floor, and population laws (with a Lean 4 formalization)},
preprint, repository snapshot
September~24, 2026.
\texttt{paper/short/main.tex}.
\href{https://github.com/marcellop71/EM/blob/d04b90f5d132a452e29df91ee09ea903b811ad92/paper/short/main.tex}{github.com/marcellop71/EM, commit d04b90f}.

\bibitem{booker-variant}
A.~R. Booker, \emph{A variant of the Euclid--Mullin sequence containing
every prime}, Journal of Integer Sequences \textbf{19} (2016), Article 16.6.4.
\href{https://arxiv.org/abs/1605.08929v1}{arXiv:1605.08929v1}.

\bibitem{booker-second}
A.~R. Booker, \emph{On Mullin's second sequence of primes},
Integers \textbf{12A} (2012), Article A4.
\href{https://arxiv.org/abs/1107.3318v2}{arXiv:1107.3318v2}.

\bibitem{pollack-trevino}
P.~Pollack and E.~Trevi\~no, \emph{The primes that Euclid forgot},
American Mathematical Monthly \textbf{121} (2014), 433--437.
\href{https://doi.org/10.4169/amer.math.monthly.121.05.433}{doi:10.4169/amer.math.monthly.121.05.433}.

\bibitem{lean}
L.~de Moura and S.~Ullrich, \emph{The Lean 4 theorem prover and programming
language}, in Automated Deduction---CADE 28, Lecture Notes in Computer
Science \textbf{12699} (2021), 625--635.
\href{https://doi.org/10.1007/978-3-030-79876-5_37}{doi:10.1007/978-3-030-79876-5\_37}.

\bibitem{mathlib}
The mathlib Community, \emph{The Lean mathematical library},
Proceedings of the 9th ACM SIGPLAN International Conference on Certified
Programs and Proofs (2020), 367--381.
\href{https://doi.org/10.1145/3372885.3373824}{doi:10.1145/3372885.3373824}.
\end{thebibliography}
\end{document}
